% ECONOMICS_PROBLEM: {"id":"G28-66","title":"Credible bank resolution","jel_code":"G28","source_id":"OP-0327"}
% FIRST_POSTED: 2026-10-11
% ENTRY_KIND: research_attempt
% ATTEMPT_STATUS: solved
% SUBMISSION: {"kind":"research_attempt","category":"economics","problem_id":"G28-66","model":{"provider":"OpenAI","version":"GPT-6 based Codex","version_uncertainty":"The system identifies GPT-6; the exact backend variant is not exposed.","reasoning":"Not exposed","generated_on":"2026-10-11"},"other_models":"Existing same-session Codex reviewers inherited model settings; no different model was knowingly selected.","tools":"Primary-source web inspection, local file and Git inspection, lightweight text/hash checks, and separate remote finite exact-fraction corroboration. No executed general CAD solver, numerical optimization, Lean proof or formal proof certificate. Release verification is recorded separately.","target":"Original G28-66 fixed-banking-system welfare-optimal credible resolution agenda, including actual clearing and runs, all declared failure states, commitment, and pricing identification.","scope":"Full fixed-system answer for the explicitly specified institution with an entire continuous capital/action domain and legal lotteries, all stress probabilities and systemic probabilities. General finite-history semialgebraic representation handles exact credibility, equilibrium/pricing, endpoint flags and epsilon-optimal rules. This is not a universal theorem for arbitrary bank laws or infinite strategic economies.","human_contributions":"Ji Ho Bae supplied the task and constraints and is the submitter; no human mathematical derivation or independent mathematical verification was reported.","verification":"Complete modules and the whole manuscript received separate same-session model-assisted reviews. The recorded specification clarifications are incorporated; exact-version audit outcomes are supplied with the release. This is not independent human verification.","reproducibility":"Canonical and both manifest-listed priors at b1f633121074e16579fa7fd05314e5f504d29bae were read. Contribution prompt A guided research and prompt B guides final source audit. Reviewed modules, source hashes and actual release-check outcomes are recorded separately."}
% FULL_SOLUTION_SCOPE: {"target":"Entire G28-66 research task for the fixed banking institution specified here, including all legal resolution deviations, both institutional variants and pricing versus sequential credibility.","result":"Priority-clearing theorems; exact legal/SE/pricing design theorem; sharp state-price theorem; complete two-bank authority and run solution; global all-probability capital algorithm and commitment classification; lottery reduction; full-continuum prices and observationally identical credibility counterexample.","coverage":"All K,p,rho in [0,1], every isolated/joint and run/wait history including zero-probability histories, all feasible continuous authority actions and admitted lotteries, initial and stress budgets, all claim seniorities, actual interbank losses, commitment/discretion, sharp identification and nonidentification. General represented institutions include infeasibility, unbounded values and nonattainment. Public capital lotteries are independent of future Nature shocks.","dependencies":"Effective real quantifier elimination and algebraic sampling from Basu-Pollack-Roy and Basu1409.1534; finite-game Nash existence for a reproduced perturbed-agent proof of sequential-equilibrium existence; Kreps-Wilson consistency definition. Clearing, economic optimization, lottery and identification arguments are proved here.","human_verification":"None reported. Same-session model-assisted review is disclosed separately.","unverified":"No independent human or formal verification, executed general CAD solver, empirical replication or exhaustive worldwide novelty certification. Exact-source model reviews and executed release checks are reported separately with their limited scope.","novelty":"Complete canonical and both prior-attempt comparison plus targeted primary-source inspection. Established clearing context and SE/QE tools and prior quadratic exposure-cost results are credited. No worldwide research-priority claim."}
\documentclass[11pt,a4paper]{article}
\usepackage[margin=25mm]{geometry}
\usepackage{amsmath,amssymb,amsthm}
\usepackage[hidelinks]{hyperref}
\hypersetup{pdftitle={Credible Bank Resolution: Clearing, Runs, and Optimal Instruments},pdfauthor={OpenAI GPT-6 based Codex}}
\setlength{\parindent}{0pt}
\setlength{\parskip}{4pt}
\setlength{\emergencystretch}{2em}
\allowdisplaybreaks[2]
\numberwithin{equation}{section}
\newtheorem{theorem}{Theorem}[section]
\newtheorem{lemma}[theorem]{Lemma}
\newtheorem{proposition}[theorem]{Proposition}
\newtheorem{corollary}[theorem]{Corollary}
\theoremstyle{definition}
\newtheorem{definition}[theorem]{Definition}
\newtheorem{remark}[theorem]{Remark}
\newcommand{\R}{\mathbb R}
\newcommand{\E}{\mathbb E}
\newcommand{\Prb}{\mathbb P}
\newcommand{\Cred}{\operatorname{Cred}}
\DeclareMathOperator{\supp}{supp}
\DeclareMathOperator{\spanlin}{span}
\title{Credible Bank Resolution:\\Clearing, Runs, and Optimal Instruments}
\author{OpenAI GPT-6 based Codex\thanks{The exact backend variant and reasoning setting are not exposed. Ji Ho Bae supplied the task and constraints and is the submitter; no independent human mathematical verification was reported.}}
\date{11 October 2026}
\begin{document}
\maketitle
\begin{center}
\small\textbf{G28-66 \quad JEL G28 \quad Source OP-0327}\\
Full original fixed-system proof claim.\\
Model-assisted source review and release verification are recorded separately.
\end{center}
\begin{abstract}
We solve a fixed banking institution's complete credible-resolution design
problem with actual interbank payments, creditor runs, insured deposits,
senior and designated debt, contractual AT1 losses, equity, bridge operations
and temporary funding. Every legal continuous authority action is compared
at every isolated, joint and off-path failure history. Paid-in capital varies
over an interval, and stress and systemic probabilities each vary over the
whole unit interval. An explicit finite candidate rule gives the global
discretionary optimum; the commitment optimum has an exact classification,
including flat boundary cases. Admitted lotteries cannot improve these
values. A particular stress probability yields a systemic-probability
threshold of \(15/49\) and a strict commitment gap. We derive claim prices
for the entire capital continuum, sharp observation-compatible sets, and
two economies with identical realized prices and behavior but opposite
credibility judgments for the same announcement. A general priority-clearing
and real-algebraic framework treats complete legal domains, equilibrium
pricing, empty feasibility, unbounded values and unattained suprema.
\end{abstract}

\section{Original target, economic scope, and context}\label{sec:target}

The original catalogue problem~\cite{catalogue}, entitled \emph{Credible
bank resolution}, fixes a banking system with insured deposits, unsecured
claims of stated seniorities, equity and cross-bank exposures. A rule
\(a\) specifies loss allocation, recapitalization, transfer of critical
functions and temporary funding. Investors price claims before stress,
anticipating the rule and actual authority deviations. The target is
\begin{equation}\label{eq:original}
\begin{split}
A_{\rm cred}
 &=\{a\in A:V^A(a(s);s,a)\ge V^A(d;s,a)
                \text{ for all declared }s,\ d\in D(s)\},\\
W_*&=\sup_{a\in A_{\rm cred}}W(a).
\end{split}
\end{equation}
The feasible set must follow the insolvency law, loss-absorbing claims and
operational constraints. Welfare includes household surplus, payments and
credit services, fiscal resources and spillovers. Runs and other banks'
write-down losses must be modeled. Isolated and joint failures must both
be assessed. Price-implied losses must be distinguished from sequential
policy optimality, under explicit pricing and recovery restrictions.

The two named variants are an irrevocably enforceable loss allocation and
a discretionary authority subject to all-state credibility. For a finite
optimal value the requested output is a feasible \(\varepsilon\)-optimal
rule for every \(\varepsilon>0\); exact attainment needs proof. Empty
feasibility and unbounded values are separate outcomes. The source's common
conventions also require a declared equilibrium selection, complete budgets
and clearing, common primitives across observations, and honest sharpness
and endpoint claims. These conventions are used throughout.

This is a model-construction research agenda, not a numerical banking
economy specified by the source. We therefore give both a general
represented-institution theorem and a fully specified institution in which
every named economic margin above is active and the entire rule domain is
solved. Its legal domain is fixed before outcomes and is not a finite menu
chosen after locating an optimum. The full-solution claim concerns that
fixed-system task. It does not assert a theorem for every possible banking
law, arbitrary non-effective utility, hidden-capital mechanism, or infinite
strategic horizon.

The catalogue attributes the empirical motivation to Di Stefano, Lengwiler
and Rishabh~\cite{bis}, and the broader policy agenda to Tucker~\cite{tucker}.
Neither source is asserted to prove the editorial welfare-optimal
credibility problem. The two earlier attempts~\cite{priorone,priortwo}
already establish a one-bank quadratic spillover optimum and a multibank
positive-definite exposure-cost characterization with observable-risk
feedback. Their exposure costs do not determine payment clearing or runs,
and their credible rule is a singleton at fixed primitives. Those results
remain credited. The contribution here is the joint solution of actual
clearing, runs, a nontrivial credible instrument continuum, commitment,
and the associated identification questions.

\section{Actual priority clearing}\label{sec:clearing}

There are finitely many banks \(i\) and classes \(r=1,\ldots,R_i\), ordered
from most to least senior. At a fixed history and proposed resolution
action, remaining face amounts are \(b_{ir}\ge0\). Protected deposits cannot
be written down. Contractual AT1 absorption, designated resolution debt
and senior debt have their specified distinct treatments. Write
\(\pi_{ir,j}\ge0\) for the fraction of class \((i,r)\) held by bank \(j\),
with \(\sum_j\pi_{ir,j}\le1\); the residual is externally held.
Outside assets and authorized bank cash are \(e_i\ge0\).

For a payment vector \(v\) in \(\prod_{ir}[0,b_{ir}]\), define
\begin{equation}\label{eq:waterfall}
\begin{aligned}
A_i(v)&=e_i+\sum_{j,r}\pi_{jr,i}v_{jr},\\
\Phi_{ir}(v)&=\min\left\{b_{ir},
                   \left[A_i(v)-\sum_{s<r}b_{is}\right]_+\right\}.
\end{aligned}
\end{equation}
A clearing vector satisfies \(v=\Phi(v)\), and equity is
\([A_i(v)-\sum_r b_{ir}]_+\). Payments within a class are proportional to
ownership. Deposit-insurance transfers to external depositors are recorded
outside this bank-cash equation. A payment legally entering another bank's
assets must instead be included explicitly. The classical network-clearing
context is Eisenberg--Noe~\cite{en}; the following proofs include priority
classes and are supplied directly.

\begin{theorem}[Unique clearing and certified errors]\label{thm:clearing}
Suppose \(\sum_j\pi_{ir,j}\le1-\eta\) for every class, for some \(\eta>0\).
There is exactly one clearing vector \(v^*\). Simultaneous iteration
converges from every vector in the payment box, with
\[
 \|v^t-v^*\|_1\le(1-\eta)^t\|v^0-v^*\|_1,\qquad
 \|v-v^*\|_1\le\eta^{-1}\|v-\Phi(v)\|_1.
\]
For fixed faces and ownership fractions,
\[
 \|v^*(e)-v^*(\widetilde e)\|_1
       \le \eta^{-1}\|e-\widetilde e\|_1.
\]
\end{theorem}
\begin{proof}
For one bank, let \(w_i(a)\) be its waterfall vector as a function of
resources \(a\ge0\). Each coordinate is nondecreasing and
\(\sum_r w_{ir}(a)=\min\{a,\sum_r b_{ir}\}\). Therefore
\[
\|w_i(a)-w_i(a')\|_1
 =\left|\min\{a,\sum_r b_{ir}\}-\min\{a',\sum_r b_{ir}\}\right|
 \le |a-a'|.
\]
Consequently
\[
\begin{split}
\|\Phi(v)-\Phi(v')\|_1
&\le\sum_i\left|\sum_{j,r}\pi_{jr,i}(v_{jr}-v'_{jr})\right|\\
&\le\sum_{j,r}\left(\sum_i\pi_{jr,i}\right)|v_{jr}-v'_{jr}|
 \le(1-\eta)\|v-v'\|_1.
\end{split}
\]
The map preserves the closed payment box. Its successive iterate
differences form a summable geometric bound, so iterates converge to a
fixed point by continuity. Two fixed points have distance at most
\(1-\eta\) times their distance and hence coincide. Comparing any iterate
with that fixed point gives the first bound. The triangle inequality gives
\(\|v-v^*\|_1\le\|v-\Phi(v)\|_1+(1-\eta)\|v-v^*\|_1\),
which proves the residual bound. Applying the same inequalities to maps
with different \(e\), and moving the contraction term to the left, proves
the final assertion.
\end{proof}

If fractions also change while faces stay fixed, the asset-difference
bound before division by \(\eta\) acquires the additional term
\[
\sum_{j,r}\widetilde v^*_{jr}
                  \sum_i|\pi_{jr,i}-\widetilde\pi_{jr,i}|.
\]
This follows by inserting the second fixed point into the first map.
The first map's leakage bound suffices. Varying faces is not covered
without extra terms.

\begin{theorem}[Clearing without leakage]\label{thm:extremeclearing}
Under merely \(\sum_j\pi_{ir,j}\le1\), least and greatest clearing vectors
exist. They have effective semialgebraic graphs for represented
semialgebraic parameters. Either may be a declared clearing convention;
uniqueness and the preceding geometric bounds are not asserted.
\end{theorem}
\begin{proof}
The map \(\Phi\) is continuous and coordinatewise nondecreasing. Iteration
from the full face vector \(b\) decreases and is bounded below by zero.
Coordinate limits exist and form a fixed point \(v^+\). Every fixed point
\(z\) satisfies \(z\le b\), hence \(z\le\Phi^t(b)\) for every \(t\), and
therefore \(z\le v^+\). Iteration from zero similarly gives the least
fixed point \(v^-\).

Let \(C(v)\) denote the box constraints and fixed-point equations
\eqref{eq:waterfall}. The greatest graph is
\[
 C(v)\ \wedge\ \forall z\,\bigl(C(z)\Longrightarrow z\le v\bigr);
\]
the least graph reverses the final inequality. The min and positive-part
operations have finite polynomial descriptions, so real quantifier
elimination gives the parameter assertion.

There is also a direct exact algorithm at rational input. Each coordinate
of \(\Phi\) is zero, its full face, or resources minus senior faces.
Enumerate all branch patterns and solve the resulting bounded linear
polyhedra. Their union is exactly the clearing set, including boundary
overlap. In the unique case every feasible branch gives the same vector.
In general maximize or minimize \(\sum_{ir}v_{ir}\) over every nonempty
branch and compare its optima. Since \(v^+\) dominates every clearing,
it is the unique maximizer of this sum; \(v^-\) is the unique minimizer.
Thus the procedure is finite even if monotone iteration never terminates
in finitely many steps.
\end{proof}

Selecting greatest or least clearing is an economic convention, not a
uniqueness conclusion. An all-clearing welfare criterion must retain the
whole correspondence \(C\) in its quantifiers. The concrete institution
below has leakage and direct acyclic clearing, so this selection issue
does not arise there.

\section{Complete legal domains, incentives, and exact design}\label{sec:general}

\subsection{The represented institution}
The following are input assumptions for a general theorem, not claims
about every real resolution regime. Coefficients are rational or explicitly
represented real algebraic numbers. All stated sets and function graphs
are given by finite real polynomial formulas.

A finite perfect-recall prefix game includes any specified bank financing
or risk decisions, Nature signals and creditor withdrawal actions. Each
terminal prefix history reaches one authority decision and then the
specified passive settlement. The authority's public history contains all
its payoff-relevant information. If unresolved private uncertainty is
present, its conditional beliefs must instead be explicit in the model;
they cannot silently be replaced by a realized public state.

The legal infrastructure parameter \(k\) ranges over a supplied
semialgebraic set. A finite declared list \(\mathcal H\) of failure/run
histories is fixed independently of their equilibrium probabilities.
At each \(h\), the entire nonempty legal action set is \(D_h(k)\).
It may be a finite union of components for discrete choices and full
continuous intervals for cash, transfers and write-downs. A complete rule
is \(a=(k,(x_h)_{h\in\mathcal H})\) subject to every initial-finance,
operational and legal constraint; their conjunction defines \(A\).
There is no finite action grid substituted for \(D_h(k)\).

Post-resolution clearing is the unique vector of
Theorem~\ref{thm:clearing}, an explicitly selected extreme vector of
Theorem~\ref{thm:extremeclearing}, or another proved represented convention.
Denote it by \(v^*(k,h,x)\). All passive recoveries, temporary funding
repayments, resources and claimant payoffs are specified. The authority
value \(V_h(k,x,v^*)\) and private payoffs are finite semialgebraic
functions. At each fixed announced rule they are uniformly bounded on the
admitted equilibrium/price domain, but no common bound over all rules
is required. A later strategic authority stage would require its own
continuation constraints; terminal settlement is an explicit assumption.

Define
\begin{equation}\label{eq:cred}
\begin{aligned}
\Cred(a)\ \Longleftrightarrow&
\ \text{for every }h\in\mathcal H\text{ and }d\in D_h(k),\\
&V_h(k,x_h,v^*(k,h,x_h))\ge V_h(k,d,v^*(k,h,d)).
\end{aligned}
\end{equation}
This recomputes clearing on the deviation's own balance sheet. Using the
announced rule's clearing for a different haircut would be incorrect.
Since there is one terminal authority decision, \eqref{eq:cred} compares
its complete continuation deviations and is precisely the original
credibility condition in this institution. Its graph is semialgebraic by
substituting the clearing and payoff graphs and eliminating the deviation
variables. No differentiability or first-order sufficiency is used.

If legally admitted rules randomize, their full representation or a proved
reduction is needed. Deterministic rules cannot silently replace them.
Section~\ref{sec:random} proves such a reduction for all the lotteries
admitted in our concrete institution.

\subsection{A genuine sequential-equilibrium graph}
For a fixed rule, Nature law and fixed prices, prune only histories
impossible because of Nature. Do not delete a legally available strategic
action because its equilibrium probability is zero. Write \(\sigma\) for
behavioral probabilities and \(\mu_I\) for beliefs over the nodes of an
information set \(I\). Let \(R_h(\sigma)\) be the product of chance and
behavior probabilities reaching node \(h\), and
\(R_I(\sigma)=\sum_{h\in I}R_h(\sigma)\).

Consistency in the sense of Kreps--Wilson~\cite{kw} is exactly the formula
\begin{equation}\label{eq:consistency}
\begin{gathered}
\forall \epsilon>0\ \exists \sigma',\mu':\
\sigma'\text{ completely mixed},\quad
\|\sigma'-\sigma\|_\infty<\epsilon,\quad
\|\mu'-\mu\|_\infty<\epsilon,\\
R_I(\sigma')>0,\qquad
\mu'_I(h)R_I(\sigma')=R_h(\sigma')
               \quad(h\in I,\ \text{every }I).
\end{gathered}
\end{equation}
All probability vectors belong to their simplices. The same \(\sigma'\)
witnesses every information set, and the Nature parameters are held fixed
inside this formula. One implication is immediate from a consistent
sequence; conversely witnesses at \(\epsilon=1/n\) form the required
sequence. Full mixing gives positive \(R_I\) precisely on the pruned game.
When Nature probabilities depend on parameters, use their finitely many
support strata first. This pruning never deletes a declared authority
history from the separate quantifier in \eqref{eq:cred}.

Let \(U_i(I;\sigma,\mu_I)\) be player \(i\)'s conditional continuation
payoff. For each of its finitely many pure continuation plans \(\tau_i\),
impose
\begin{equation}\label{eq:seqrat}
U_i(I;\sigma,\mu_I)\ge
                 U_i(I;\tau_i,\sigma_{-i},\mu_I).
\end{equation}
These are finite polynomial comparisons after introducing the payoff
graphs. Perfect recall means no information set repeats along a path, so
the continuation expression is affine in each deviating behavioral
coordinate separately. Maximizing such an expression over the product
of simplices has a pure-plan maximizer, by maximizing one coordinate at
a time. Thus \eqref{eq:seqrat} tests all randomized continuation deviations.
Together \eqref{eq:consistency}--\eqref{eq:seqrat} give an exact finite real
formula for sequential assessments, not arbitrary off-path beliefs.

\begin{lemma}[Existence for the fixed finite prefix]\label{lem:se}
A finite perfect-recall prefix game with fixed terminal payoffs has a
sequential equilibrium after the stated Nature pruning.
\end{lemma}
\begin{proof}
If there are no strategic information sets, the empty assessment proves
the assertion. Otherwise proceed as follows.
Fix \(\delta>0\) smaller than the reciprocal of every information-set
action count \(m_I\). Replace each behavioral simplex by
\(\sigma_I(a)=\delta+(1-m_I\delta)z_I(a)\), with \(z_I\) in a simplex.
Treat every information set as an agent whose payoff is the original
player's ex ante payoff. This is a finite agent normal-form game: fixing
an agent's pure choice determines the displayed noisy action, and
independent mixing of agents gives exactly these behavioral probabilities.
The finite-game Nash theorem~\cite[Theorem~1]{nash} gives an agent
equilibrium.

Every remaining information set has positive reach. Perfect recall makes
its reach independent of that agent's current action and makes its
conditional node probabilities invariant to the original player's earlier
behavior once that information set is reached: the player's own past
action factors are common to its nodes and cancel. Consequently each
agent best response is a conditional best response among floor-constrained
actions. Backward replacement, starting at the last information sets,
shows that no floor-constrained continuation plan improves at any
information set. At a later information set the same cancellation of the
player's past factors preserves the conditional distribution during this
replacement. This proves the full continuation comparison, not merely
a condition on reached actions.

Take a sequence \(\delta\downarrow0\) and a convergent subsequence of the
behavior and Bayes-belief vectors, which lie in a finite product of compact
simplices. It is consistent because each entire profile is completely
mixed with the same fixed Nature law. For any pure continuation plan,
approximate it by putting mass \(\delta\) on each other action and the
remaining mass on its prescribed action. The floor-constrained comparison
just proved passes to the limit: conditional payoff expressions are
polynomials in the behavior and belief coordinates. It gives
\eqref{eq:seqrat} for every pure continuation plan at every information
set. There are finitely many such comparisons. Hence the limit is a
sequential equilibrium.
\end{proof}

On \(\Cred(a)\), an indifferent authority honors its prescribed optimal
continuation as an explicit selection convention. If the target instead
ranges over all optimal authority continuations, include all of them in
the equilibrium graph. Let \(\mathcal E(a)\) contain the prefix assessment,
all declared budgets, and competitive prices
\[
 Q_c=\E[MZ_c]
\]
for the specified common pricing kernel and actual payoff \(Z_c\).
If prices enter private incentives they enter those same payoff functions.
Lemma~\ref{lem:se} does not prove nonemptiness after imposing these extra
price-consistency equations. Nonemptiness is separately proved in the
concrete model, and is tested in the general graph.

\subsection{The global design theorem}
A selected-equilibrium target includes its represented selection graph.
Alternatively, declare robust-prefix welfare
\(W_-(a)=\inf_{e\in\mathcal E(a)}W(a,e)\) on nonempty \(\mathcal E(a)\).
The exact graph of its finite value \(w\) is
\begin{equation}\label{eq:lowergraph}
\begin{split}
&\forall e\in\mathcal E(a),\quad W(a,e)\ge w,\\
&\forall \epsilon>0\quad
 \exists e\in\mathcal E(a),\quad W(a,e)<w+\epsilon.
\end{split}
\end{equation}
The second condition matters when the infimum is not attained. This
robust convention is a declared target, not an implicit choice imposed
on the original question.

\begin{theorem}[Exact design over the entire represented domain]\label{thm:design}
For the represented institution and either a represented selected welfare
or \eqref{eq:lowergraph}, there is an exact finite algorithm for the
feasible credible set, its welfare value set and supremum. It distinguishes
empty feasibility, a finite supremum and an unbounded value. In the finite
case it decides attainment and returns an exact optimizer when attained.
For every specified positive rational or algebraic \(\epsilon\), it returns
an algebraic feasible credible rule with welfare greater than
\(W_*-\epsilon\). The corresponding existence statement holds for every
real \(\epsilon>0\). No polynomial running-time claim is made.
\end{theorem}
\begin{proof}
Conjoin \(a\in A\), \(\Cred(a)\), \(\mathcal E(a)\ne\varnothing\), and the
welfare graph. Everything is a finite real formula by
\eqref{eq:waterfall}--\eqref{eq:lowergraph}. Effective real quantifier
elimination and real-algebraic sampling~\cite{bpr,basu} apply: their
inputs are finitely many polynomial conditions with represented algebraic
coefficients; nonpolynomial semialgebraic functions have been replaced
by their graphs. The output describes the exact projected value set
\(S\subset\R\), a finite union of intervals and points with algebraic
finite endpoints.

Ordinary real-algebraic sampling also follows constructively from this
elimination operation. Project a nonempty represented graph to one
coordinate. Univariate real-root isolation describes its nonempty intervals
and isolated algebraic points; choose a rational interior point of an
interval, or an isolated algebraic point. Fix that coordinate and repeat
on the nonempty fiber. After finitely many coordinates this gives an
ordinary real-algebraic witness, not merely a point in an infinitesimal
extension.

If \(S\) is empty, feasibility fails. Otherwise its upper endpoint is
algebraic or \(+\infty\). For a finite endpoint \(s\), membership of \(s\)
decides attainment, and sampling the original graph at that value returns
an optimizing rule and any required witnesses. If it is not attained,
the defining property of the supremum makes the graph with
\(W>s-\epsilon\) nonempty. Algebraic sampling again gives the legal rule,
not only a numerical lower bound on its value. For an arbitrary positive
real tolerance, choose a smaller positive rational tolerance to obtain
existence. If the value is unbounded, this is reported separately.
\end{proof}

These real-algebraic tools are established results, not a novelty claim.
Their content used here is effective elimination of finite real quantifiers
and extraction of algebraic sample points; see
\cite[Theorem~2.27]{basu} and the sampling construction above. The many continuation plans and
quantifier blocks may be expensive. This manuscript does not claim to
have executed a general CAD or quantifier-elimination computation.

\section{What common-kernel claim prices identify}\label{sec:generalid}

Fix a complete finite settlement-state list \(s=1,\ldots,L\), including
normal repayment. If only stress states are retained, explicitly subtract
the known normal-state price contribution. Let \(R_{sc}\) be the known
actual payoff of claim \(c\), incorporating runs and clearing.
For this lemma the physical probabilities \(\pi_s>0\) are known and a
common positive pricing kernel obeys known finite bounds
\(0<\ell_s\le M_s\le u_s\). Put \(q_s=\pi_sM_s\) and define
\begin{equation}\label{eq:pricepoly}
\mathcal D(Q)=
\{q:R^\top q=Q,\quad \pi_s\ell_s\le q_s\le\pi_su_s\ \forall s\}.
\end{equation}
A safe unit-payoff claim, if observed, is a column of \(R\) and fixes
\(\sum_s q_s\).

\begin{theorem}[Sharp state-price information]\label{thm:pricepoly}
\(\mathcal D(Q)\) is exactly the identified state-price set. It is empty
precisely for incompatible prices and restrictions. If nonempty, the
sharp bounds on \(t^\top q\) are its attained linear-program extrema.
That target is point identified exactly when
\[
t\perp\spanlin(\mathcal D(Q)-\mathcal D(Q)).
\]
If a point of \(\mathcal D(Q)\) is strictly inside every box bound, this
is equivalent to \(t\) belonging to the row space of \(R^\top\).
\end{theorem}
\begin{proof}
Every admitted kernel gives \eqref{eq:pricepoly}. Conversely a retained
\(q\) gives \(M_s=q_s/\pi_s\), satisfying the bounds and every price
equation. This proves sharpness. The set is closed and bounded, so the
linear extrema are attained. A linear functional is constant exactly
when it annihilates every feasible difference. At a strictly interior
point \(q_0\), every \(v\in\ker R^\top\) has \(q_0+tv\) feasible for all
sufficiently small \(|t|\). Thus the difference span is the kernel, whose
orthogonal complement is the row space.
\end{proof}

Binding kernel bounds can identify a target without the row-space
criterion, so the interior qualification cannot be dropped. If physical
probabilities are unknown and null states are allowed, the sharp joint
set instead uses
\[
\pi\ge0,\quad \sum_s\pi_s=1,\quad
R^\top q=Q,\quad \ell_s\pi_s\le q_s\le u_s\pi_s.
\]
On a null state these inequalities force \(q_s=0\), and any in-bound
\(M_s\) can be assigned there; on a positive state use \(M_s=q_s/\pi_s\).
This also proves the converse. If strictly positive support is required,
replace \(\pi\ge0\) by \(\pi>0\); compactness may fail and endpoints must
be reported as infima/suprema with their actual attainment flags.

Even complete state prices do not separate an unknown physical law from
its kernel or reveal the authority's ranking of unused legal deviations.
For a represented structural class, impose all observed prices, run and
payment probabilities, and parameter restrictions in one common graph.
Projection gives sharp joint sets for the credibility indicator, optimal
instruments and welfare. An announced rule and actual discretionary
actions are separate variables: their equality must be tested, not
assumed in order to infer credibility. Section~\ref{sec:prices} supplies
both a constructive sharp example and an explicit nonidentification pair.

\section{A fully specified banking institution}\label{sec:institution}

\subsection{Balance sheets, timing, and funded capital}
Fix banks \(A,B\) and put \(\gamma=1/2\). Bank \(A\) has insured deposits
of face \(1\), protected senior unsecured debt of face \(1/2\), designated
resolution debt of face \(1\), contractual AT1 of face \(1/4\), and old
equity of \(1/4\). In normal conditions its assets \(3\) pay all debt and
leave old equity \(1/4\). In stress its pre-resolution assets are \(3/2\):
old equity is wiped and AT1 contractually absorbs \(1/4\). The remaining
debt faces total \(5/2\), leaving recapitalization shortfall \(1\).
AT1 absorption is not a discretionary resolution haircut.

Bank \(B\) holds fraction \(\gamma\) of \(A\)'s designated debt; its other
half and all other claims are held outside the bank network. An \(A\)
haircut \(x\) therefore reduces \(B\)'s actual payment receipt by
\(\gamma x\). This is an interbank asset, not an assumed spillover cost.

At date zero the designer publicly chooses paid-in capital \(K\in[0,1]\).
It contributes cash \(\gamma K\) to \(B\) before stress and has a real
capital-raising/locking friction \(K/2\). Cash principal is a financial
contribution with residual ownership, not another social loss. Adequately
funded outside owners supply it through a legally binding date-zero
capital assessment. They retain the stated residual ownership; voluntary
participation or withdrawal is not an additional unmodeled decision.
The contribution is not a promise of future rescue and is locked before
the terminal authority acts. Any public lottery over
\(K\) is drawn before and independently of the subsequent Nature shocks,
with no preliminary signal about them. In particular, the conditional
shock law below remains \((p,\rho)\) given its realization.

Bank \(B\) has insured deposits \(\gamma\) and designated wholesale debt
\(\gamma\). In a stress state \(H\), its external assets are
\(\gamma(1+K-H)\). With the actual incoming \(A\) payment they become
\begin{equation}\label{eq:Bassets}
 \gamma(1+K-H)+\gamma(1-x)=\gamma(2+K-H-x).
\end{equation}
The declared states are \(H=0\), an isolated initial loss at \(A\), and
\(H=1/2\), an additional external-asset loss at \(B\). At \(K=0\) the
latter makes both banks initially insolvent. Capital and incoming
write-downs endogenously change subsequent shortfalls. In normal
conditions \(H=x=0\), \(B\)'s faces are fully paid and equity is
\(\gamma K\).

Nature generates stress with probability \(p\in[0,1]\); conditional on
stress, \(H=1/2\) has probability \(\rho\in[0,1]\). Normal claims are
fully paid. Both hypothetical stress states and both creditor histories
remain in the declared credibility list even at probability endpoints.
The probability square is an exogenous primitive family, independent
of capital choices. At \(p=0\), \(\rho\) labels counterfactual stress
composition rather than an observable positive-probability conditional law.

\subsection{Runs, temporary liquidity, and the entire legal domain}
After observing \(K,H\) but before resolution, the outside wholesale
creditor chooses \emph{wait} or \emph{run}. Waiting retains its designated
claim. A run exercises early redemption: a standing central-bank facility
lends \(\gamma\) to \(B\), which pays \(\gamma\) to a settlement channel.
The creditor receives \(\gamma/2\); the remaining \(\gamma/2\) is a real
processing/fire-sale cost. Bank cash changes by
\(+\gamma-\gamma=0\), and the wholesale liability is replaced by a
protected central-bank liability \(\gamma\). Its pre-resolution assets
\eqref{eq:Bassets} are unchanged. The advance is repaid after
recapitalization; it is not counted again as an unrecovered fiscal loss.
The standing advance is a binding right of this fixed institution under
both variants. Withholding it is not a legal terminal deviation.

A run also incurs real critical-service interruption \(7/20\). Insured
deposits and \(A\)'s senior unsecured debt are protected. Required payment
functions continue through the statutory bridge operation; liquidation
without those functions is outside the stated operational law.
For waiting, the authority's \emph{entire} legal action set is
\begin{equation}\label{eq:legal}
 0\le x\le1,\qquad
 0\le y\le\min\{1,[H+x-K]_+\}.
\end{equation}
Here \(x\) is \(A\)'s designated haircut and \(\gamma y\) is \(B\)'s.
A solvent bank's equity cannot remain while its designated debt is
impaired. All old insolvent equity is wiped. The minimum required public
bank cash is
\begin{equation}\label{eq:cash}
 b_A=1-x,\qquad b_B=\gamma\bigl([H+x-K]_+-y\bigr).
\end{equation}
Excess cash gifts beyond restoration of remaining debt faces are legally
disallowed; they are not an omitted authority action.

If \(H+x-K>0\), \(B\)'s assets plus \(b_B\) are
\(\gamma(2-y)\), exactly paying insured deposits \(\gamma\) and wholesale
debt \(\gamma(1-y)\). Otherwise \(b_B=y=0\), all faces are paid, and
equity is \(\gamma(K-H-x)\). Bank \(A\)'s estate is
\(3/2+b_A=5/2-x\), exactly paying its protected \(3/2\) and designated
\(1-x\). Thus the actual interbank payment and both priority waterfalls
clear. These identities also prove clearing directly without iteration.

After a run \(y=0\), because the replacement central-bank claim is
protected, and \emph{every} \(x\in[0,1]\) remains a legal deviation.
Formula~\eqref{eq:cash} restores the needed estate in that branch too.
Bridge operating purchases cost \(x/4\). A stress-date cash capacity of
\(2\) covers all legal actions: \(b_A+b_B\le5/4\), the advance is at most
\(1/2\), and operating purchases are at most \(1/4\). The advance counts
toward peak liquidity, although it is later repaid. Operating inputs are
purchased from a prefunded reserve at their real opportunity cost.
Date-zero capital cash and its raising friction are separate date-zero
flows. Outside household/taxpayer cash may be fixed at \(10\), so none
of these finite funding requirements relies on an unavailable resource.

The complete announced rule specifies the action in
\eqref{eq:legal} at both waiting states and an \(x\) after each run,
as well as \(K\). Recapitalization and bridge/temporary-funding operations
are the legally specified functions just given. Thus no feasible
resolution action has been discarded by representing the rule this way.
Arbitrary legal action lotteries are addressed in Section~\ref{sec:random}.

\subsection{Credit production, private payoffs, and welfare}
Wholesale receipts fund a separate working-capital sector through a
pass-through par lending facility. With normalized funds \(z\in[0,1]\),
entrepreneurs choose \(0\le q\le z\), use working capital \(\gamma q\),
and produce gross resources
\(\gamma q+q-q^2/2\). The principal \(\gamma q\) is repaid at par.
Net real surplus
\[
 F(q)=q-q^2/2
\]
is maximized at \(q=z\). Entrepreneurs own that surplus; the creditor's
financial payoff remains its cash repayment, subsequently lent and
returned at par. With unit discounting, this pass-through use does not
silently add entrepreneurial surplus to the debt price or run incentive.
Normal \(z=1\) gives \(F(1)=1/2\). Waiting gives \(z=1-y\), so the
production loss is \(y^2/2\). Running gives \(z=1/2\), so it is \(1/8\).

The social numeraire is the real consumption good, with common unit
weights on its owners' resources. Public contributions \(b_A,b_B\) have
real fiscal distortions \(b_A\) and \(2b_B\), respectively; principal
transfers are not counted a second time. Prefunded bridge operating inputs
have total real cost \(x/4\). Conditional on waiting, the authority
minimizes precisely
\begin{equation}\label{eq:loss}
 L_H(x,y;K)
 =1-x+[H+x-K]_+-y+x/4+y^2/2.
\end{equation}
After running its action-dependent loss is \eqref{eq:loss} at \(y=0\).
The additional loss fixed by that branch is
\begin{equation}\label{eq:runloss}
 R=\gamma/2+7/20+1/8=29/40.
\end{equation}
This includes processing, service interruption and the resulting
production loss. It cannot be reversed by the authority's \(x\) choice
but is included in ex ante welfare.

All other shock losses and operating surplus are common across designs
at fixed \((p,\rho)\). Their baseline \(W_0(p,\rho)\) may depend on the
shock law. Thus
\[
 W(K;p,\rho)=W_0(p,\rho)-J(K;p,\rho).
\]
The baseline cancels in matched policy/institution comparisons at the
same primitives, not automatically across distinct laws.

The creditor receives \(\gamma(1-y)\) after waiting and \(\gamma/2\)
after running. It waits when the anticipated waiting haircut is at most
\(1/2\). At indifference select waiting. This is an explicit
selected-equilibrium convention, not a claim that all indifferent
creditors behave identically. Existing claims trade between adequately
funded outside households; price is their valuation, not an additional
unmodeled issuer-financing constraint. Initial capital funding was
separately supplied above.

\section{Every authority deviation and the endogenous run}\label{sec:authority}

\begin{theorem}[The complete discretionary continuation]\label{thm:authority}
For every \(K\in[0,1]\) and \(H\in\{0,1/2\}\), the unique authority
optimum after waiting is
\begin{equation}\label{eq:waitopt}
 x_W=\min\{1,K-H+3/4\},\qquad
 y_W=\min\{3/4,1+H-K\}.
\end{equation}
After a run the unique optimum is
\begin{equation}\label{eq:runopt}
 x_R=[K-H]_+,\qquad y_R=0.
\end{equation}
With the stated creditor selection, the run rule is
\[
 H=0:\ K<1/2;\qquad H=1/2:\ K<1.
\]
The rule is credible at both isolated and joint states and at both
creditor histories, independently of their probabilities.
\end{theorem}
\begin{proof}
Fix a feasible \(y\). If \(y>0\), legality requires \(x\ge K-H+y\).
Above the resource kink the part of the loss varying with \(x\) has
slope \(1/4>0\). If \(y=0\) and \(K>H\), it has slope \(-3/4\) up to
\(x=K-H\) and \(1/4\) thereafter. If \(K\le H\), the kink is at or
below zero. In all cases the unique minimizing \(x\) is
\[
 x^*(y)=[K-H+y]_+,
 \qquad 0\le y\le\min\{1,1+H-K\}.
\]
Substitution into \eqref{eq:loss} gives the strictly convex function
\begin{equation}\label{eq:reduced}
 f_H(y;K)=1+H-K-y+\frac{[K-H+y]_+}{4}+\frac{y^2}{2}.
\end{equation}
Before the kink its derivative is \(y-1\); after it the derivative is
\(y-3/4\). Since \(H-K\le1/2\), the derivative is still negative
through any nonnegative kink. The minimum is therefore \(y=3/4\)
unless the upper endpoint \(1+H-K\) intervenes. This proves
\eqref{eq:waitopt}; the lower \(x\) bound is automatic because
\(K-H+3/4\ge1/4\). Strict convexity in \(y\), and uniqueness of \(x^*(y)\),
prove uniqueness over the entire legal set.

After running \(y=0\); the same two slopes give \eqref{eq:runopt},
whose upper bound is valid since \(K-H\le1\). The fixed loss \(R\)
does not affect this minimization. Finally the cash comparison says
run exactly when \(y_W>1/2\). For \(H=0\), \eqref{eq:waitopt} crosses
that level at \(K=1/2\); for \(H=1/2\) it crosses at \(K=1\).
The equalities use the declared wait selection. This proves all
continuation comparisons, including off-path histories.
\end{proof}

The proof gives feasible behavior and optimal continuations directly,
so the concrete equilibrium exists. Before stress, competitive
risk-neutral valuation assigns the expected actual repayments derived
below. Those prices do not feed an additional financing decision, so
there is no unsolved price-consistency fixed point.

\section{The global credible capital optimum}\label{sec:capital}

Write \(S_0(K)\) and \(S_J(K)\) for the conditional real stress losses,
including \(R\) when a run occurs. Direct evaluation of
\eqref{eq:waitopt}--\eqref{eq:runopt} gives
\begin{equation}\label{eq:conditional}
\begin{aligned}
S_0(K)&=
\begin{cases}
69/40-3K/4,&0\le K<1/2,\\
1/4+(1-K)^2/2,&1/2\le K\le1,
\end{cases}\\
S_J(K)&=
\begin{cases}
89/40-K,&0\le K<1/2,\\
21/10-3K/4,&1/2\le K<1,\\
3/8,&K=1.
\end{cases}
\end{aligned}
\end{equation}
For example, in the single run state \(x_R=K\), so the variable loss
is \(1-3K/4\); adding \(29/40\) gives the first line. In the joint
run state \(x_R=0\) below \(1/2\), giving \(3/2-K+29/40\);
above \(1/2\), \(x_R=K-1/2\), giving
\(3/2-K+(K-1/2)/4+29/40\). In a single waiting state
\(x_W=1,y_W=1-K\); in the joint waiting state at \(K=1\),
\(x_W=1,y_W=1/2\). These give all remaining lines and prove
\eqref{eq:conditional} without a numerical optimization.

The ex ante loss is
\begin{equation}\label{eq:Jdef}
 J(K;p,\rho)=K/2+p\bigl((1-\rho)S_0(K)+\rho S_J(K)\bigr).
\end{equation}
Equivalently,
\begin{equation}\label{eq:J}
J(K;p,\rho)=
\begin{cases}
p(69/40+\rho/2)+K[1/2-p(3+\rho)/4],
       &0\le K<1/2,\\
\begin{aligned}
p(3/4+27\rho/20)&+K(1/2-p+p\rho/4)\\
                &+p(1-\rho)K^2/2
\end{aligned}&1/2\le K<1,\\
1/2+p(1/4+\rho/8),&K=1.
\end{cases}
\end{equation}
At \(1/2\) the value drops by \(39p(1-\rho)/40\) relative to its
left limit; at \(1\) it drops by \(39p\rho/40\). These jumps may
vanish at probability endpoints, but never point upward.

\begin{theorem}[Exact design for the whole probability square]\label{thm:capital}
For every \((p,\rho)\in[0,1]^2\), define
\[
\mathcal C(p,\rho)=\{0,1/2,1\}
\]
and adjoin
\begin{equation}\label{eq:Kc}
 K_c=\frac{p(1-\rho/4)-1/2}{p(1-\rho)}
\end{equation}
exactly when \(p>0,\rho<1\) and \(1/2<K_c<1\). A least value of
\eqref{eq:J} on this finite set is the global minimum over all
\(K\in[0,1]\), and its associated rule in
Theorem~\ref{thm:authority} is an exact welfare-optimal credible rule.
\end{theorem}
\begin{proof}
The first branch is affine. Its infimum is at zero or its right limit,
and the actual value at \(1/2\) weakly improves that limit. The middle
branch is convex quadratic. When its curvature is positive, its only
possible strict interior minimum is the zero of its derivative,
\eqref{eq:Kc}; otherwise it is affine or constant. Its endpoint
possibilities are \(1/2\) and the left limit at \(1\), which the actual
value at \(1\) weakly improves. This proves the candidate rule even at
\(p=0\) or \(\rho=1\), without dividing by zero, and proves attainment
despite selected-equilibrium jumps. A minimizing candidate suffices;
the assertion does not need to enumerate every point of a possible
flat branch.
\end{proof}

\begin{corollary}[The systemic threshold at \(p=1/2\)]\label{cor:threshold}
When \(p=1/2\), the optimizer is uniquely \(K=1/2\) for \(\rho<15/49\),
uniquely \(K=1\) for \(\rho>15/49\), and exactly the two points
\(\{1/2,1\}\) at equality.
\end{corollary}
\begin{proof}
The first branch becomes \(69/80+\rho/4+K(1-\rho)/8\), and the
middle derivative is \(\rho/8+(1-\rho)K/2>0\).
Thus only \(0,1/2,1\) need be compared. Their exact differences are
\[
\begin{aligned}
J(0)-J(1/2)&=17(1-\rho)/40,\\
J(0)-J(1)&=19/80+3\rho/16,\\
J(1/2)-J(1)&=(-15+49\rho)/80.
\end{aligned}
\]
For completeness,
\(J(1/2)=7/16+27\rho/40\) and \(J(1)=5/8+\rho/16\).
The only flat first branch is \(\rho=1\), whose value is still strictly
above \(J(1)\). Hence there are no additional optimizers.
\end{proof}

\begin{corollary}[Rectangular ambiguity in incremental costs]\label{cor:robust}
At fixed \(K\), \(J\) is nondecreasing in both \(p\) and \(\rho\).
For a declared incremental-cost criterion and uncertainty rectangle
\([p_\ell,p_u]\times[\rho_\ell,\rho_u]\), its worst case is
\((p_u,\rho_u)\). Theorem~\ref{thm:capital} at that corner is an exact
minimax incremental-cost design.
\end{corollary}
\begin{proof}
Both conditional losses are nonnegative. Their difference is
\(1/2-K/4>0\) below \(1/2\),
\(27/20+K/4-K^2/2>0\) on \([1/2,1)\), and \(1/8\) at \(1\).
The middle expression is concave and positive at both endpoints, which
proves its stated sign. Formula~\eqref{eq:Jdef} proves monotonicity.
The same upper corner is a worst case for every \(K\), so minimizing
there gives the minimax rule. A claim about total welfare across laws
would also have to include \(W_0(p,\rho)\); this corollary does not
silently cancel that baseline.
\end{proof}

\section{Irrevocable commitment on the same legal domain}\label{sec:commitment}

In the commitment variant the chosen loss allocation is irrevocably
implemented. The balance sheets, protected claims, standing liquidity
facility, capital instrument and resource constraints do not change.
To induce waiting a deterministic committed rule must have \(y\le1/2\).

\begin{theorem}[Conditional commitment and its global optimum]\label{thm:commitment}
For every \(K,H\), a best committed rule induces waiting. Its statewise
actions and losses are
\begin{equation}\label{eq:commitstates}
\begin{array}{c|c|c}
\text{state and range}&(x,y)&\text{stress loss}\\ \hline
H=0,\ K\le1/2&(K+1/2,1/2)&3/4-3K/4\\
H=0,\ K\ge1/2&(1,1-K)&1/4+(1-K)^2/2\\
H=1/2,\ 0\le K\le1&(K,1/2)&9/8-3K/4 .
\end{array}
\end{equation}
The global loss is
\begin{equation}\label{eq:Jcommit}
J_{\rm com}(K;p,\rho)=
\begin{cases}
p(3/4+3\rho/8)+K(1/2-3p/4),&0\le K\le1/2,\\
\begin{aligned}
p(3/4+3\rho/8)&+K(1/2-p+p\rho/4)\\
              &+p(1-\rho)K^2/2
\end{aligned}&1/2\le K\le1 .
\end{cases}
\end{equation}
Its complete set of minimizers is
\begin{equation}\label{eq:commitclass}
\begin{cases}
\{0\},&p<2/3,\\
[0,1/2],&p=2/3,\ \rho<1,\\
[0,1],&p=2/3,\ \rho=1,\\
\{1/2\},&p>2/3,\ p\le2/(2+\rho),\\
\{1\},&p>2/3,\ p>2/(2+\rho),\ p\rho\ge2/3,\\
\{K_c\},&p>2/3,\ p>2/(2+\rho),\ p\rho<2/3 .
\end{cases}
\end{equation}
The last case always has \(p>0,\rho<1\) and \(1/2<K_c<1\).
\end{theorem}
\begin{proof}
The reduction \eqref{eq:reduced} applies to every legal waiting action.
For \(y\le1/2\) its derivative is negative on either side of the kink.
Thus the best feasible \(y\) is \(\min\{1/2,1+H-K\}\), and
\(x=[K-H+y]_+\). Substitution gives \eqref{eq:commitstates}, including
the shared boundary \(K=1/2\).

Inducing a run cannot improve this feasible waiting commitment.
For the single state below \(K=1/2\), the least run loss exceeds the
displayed waiting loss by \(39/40\). For the joint state below \(1/2\),
the difference is \(11/10-K/4\ge39/40\); at and above \(1/2\) it is
\(39/40\). These follow by subtracting \eqref{eq:commitstates} from
the minimized run loss in \eqref{eq:runopt} plus \(R\).
In the single state above \(1/2\), the unconstrained waiting optimum
already has \(y\le1/2\). Its loss is no larger than the minimum of
the \(y=0\) subset, and a run adds the positive constant \(R\).
This argument also bounds any hypothetical run at a parameter where the
wait-at-tie rule makes it impossible to induce. A nonminimizing run
allocation cannot do better.

Taking the state expectation and adding \(K/2\) proves
\eqref{eq:Jcommit}. The two formulas agree at \(1/2\).
The left derivative is \(1/2-3p/4\). The right derivative at \(1/2\)
is \(1/2-p(2+\rho)/4\), larger by \(p(1-\rho)/4\), and thereafter
increases with slope \(p(1-\rho)\). At \(1\) it equals
\(1/2-3p\rho/4\). Hence the entire function is convex.

If \(p<2/3\), all derivatives are positive and the optimum is zero.
At \(p=2/3\), the left interval is flat; the right derivative is
positive unless \(\rho=1\), when the whole interval is flat.
For \(p>2/3\), the left derivative is negative. A nonnegative right
derivative at \(1/2\) is exactly \(p\le2/(2+\rho)\), giving that
endpoint. Otherwise a nonpositive derivative at \(1\) is exactly
\(p\rho\ge2/3\), giving \(1\). In the remaining case the derivative
has its unique zero \(K_c\) strictly inside the right interval.
That case cannot have \(\rho=1\); its curvature is positive.
At endpoint derivative equalities with \(\rho<1\), positive curvature
prevents an additional flat interval. This proves every case in
\eqref{eq:commitclass}.
\end{proof}

\begin{corollary}[A strict commitment comparison]\label{cor:gap}
At \(p=1/2\), commitment has unique optimum \(K=0\), with
\[
 J_{\rm com}^*=3/8+3\rho/16.
\]
Its waiting rules are \((x,y)=(1/2,1/2)\) in the isolated state and
\((0,1/2)\) in the joint state. If \(J_{\rm cred}^*\) denotes the
discretionary minimum, then
\begin{equation}\label{eq:gap}
J_{\rm cred}^*-J_{\rm com}^*=
\begin{cases}
1/16+39\rho/80>0,&\rho\le15/49,\\
1/4-\rho/8\ge1/8,&\rho\ge15/49.
\end{cases}
\end{equation}
\end{corollary}
\begin{proof}
The first assertion is the \(p<2/3\) case of
Theorem~\ref{thm:commitment}. Evaluate \eqref{eq:Jcommit} at zero
and \eqref{eq:commitstates} at its two states. Subtract this loss
from \(7/16+27\rho/40\) or \(5/8+\rho/16\), as appropriate in
Corollary~\ref{cor:threshold}, to get \eqref{eq:gap}. The two
expressions coincide at \(15/49\). Neither committed waiting rule is
the discretionary minimum at \(K=0\), so the comparison genuinely
uses irrevocable enforcement rather than a noncredible announcement.
\end{proof}

The comparison is at identical \((p,\rho)\) and physical/legal resources.
Its strict gap is asserted at \(p=1/2\), not for every probability:
for example at \(p=0\) both institutions optimally choose \(K=0\).

\section{The whole admitted randomized rule domain}\label{sec:random}

Allow arbitrary Borel lotteries over legal authority actions, every
realization satisfying \eqref{eq:legal} and the cash/protection constraints.
Paid-in capital is public before the run choice. A date-zero lottery
over it is independent of later Nature, as specified in
Section~\ref{sec:institution}. Hidden capital or a preliminary signal
about future shocks would be a different information structure.

\begin{theorem}[Randomization does not improve the design values]\label{thm:lottery}
The deterministic optima in Theorems~\ref{thm:capital} and
\ref{thm:commitment} are also optimal over this entire randomized rule
domain. Sequentially credible authority lotteries are almost surely
degenerate at their deterministic statewise minima.
\end{theorem}
\begin{proof}
At fixed \(K,H\), Theorem~\ref{thm:authority} gives a unique global
minimum in each creditor history. The expected nonnegative loss gap
of an optimal authority lottery must be zero. Consequently that gap
is zero almost surely, which forces the unique action. This applies
also to declared histories of zero equilibrium probability.

Under commitment, first fix \(K,H\) and a lottery over waiting actions
that has not been revealed when the creditor decides. Replace \(x\)
pointwise by \(x^*(y)=[K-H+y]_+\). This is legal and weakly lowers
real cost, preserving the creditor's cash \(\gamma(1-y)\).
The possible \(y\)'s lie in the fixed interval
\([0,\min\{1,1+H-K\}]\). Their mean \(\bar y=\E y\) lies there too,
and \(\bar x=[K-H+\bar y]_+\) is legal. Convexity of
\eqref{eq:reduced} gives
\[
 f_H(\bar y;K)\le\E f_H(y;K).
\]
In this particular formula the inequality follows directly from
\(\E y^2\ge(\E y)^2\) and
\(\E[K-H+y]_+\ge[K-H+\E y]_+\).
The creditor's expected waiting cash remains
\(\gamma(1-\E y)\), so its comparison to \(\gamma/2\), including
the declared wait-at-tie choice, is preserved. Hence the best waiting
lottery cannot improve on the deterministic waiting commitment.
If running is selected, its fixed resource losses cannot be undone,
and an action lottery cannot improve the unique run minimum.
The proof of Theorem~\ref{thm:commitment} shows that inducing this
run never beats the best feasible no-run commitment at the same \(K,H\).

If a signal or action lottery is revealed before the creditor moves,
condition on that signal. Apply the same argument to the residual
unrevealed distribution; each conditional result is bounded below by
the same best deterministic cost at that \(K,H\). Taking expectations
proves the result for mixtures of revealed and unrevealed randomization.
All actions and losses are bounded, so these expectations exist.

Finally a public capital lottery averages losses at its realized
\(K\)'s and cannot be lower than their global deterministic minimum.
This last expectation uses the unchanged conditional shock parameters
\((p,\rho)\): the lottery is before and independent of Nature and
contains no signal about future shocks. Without that primitive the
displayed averaging argument would be unjustified.
\end{proof}

This reduction uses the model's cash-linear creditor payoff, the convex
reduced cost, public capital and specified timing. It does not assume
all legal randomized policies away. Nor does it prove robustness to a
different creditor equilibrium selection at indifference.

\section{All-capital prices, sharp sets, and credibility}\label{sec:prices}

\subsection{Actual repayments and the complete price functions}
Take competitive risk-neutral pricing with unit discounting. Let \(Q_A\)
be the price per unit face of \(A\)'s designated debt, and \(Q_B\) the
price of \(B\)'s entire face-\(\gamma\) wholesale claim. Entrepreneurial
surplus belongs to other owners and is not an extra cash payoff on that
claim.

\begin{proposition}[The entire discretionary price curve]\label{prop:prices}
For every \(K,p,\rho\in[0,1]\), selected discretionary prices are
\begin{equation}\label{eq:QA}
Q_A=
\begin{cases}
1-p(1-\rho)K,&0\le K<1/2,\\
1-p[(1-\rho)+\rho(K-1/2)],&1/2\le K<1,\\
1-p,&K=1,
\end{cases}
\end{equation}
and
\begin{equation}\label{eq:QB}
Q_B=
\begin{cases}
\gamma(1-p/2),&0\le K<1/2,\\
\gamma\{1-p[(1-\rho)(1-K)+\rho/2]\},&1/2\le K<1,\\
\gamma(1-p\rho/2),&K=1 .
\end{cases}
\end{equation}
The unconditional run probability is
\begin{equation}\label{eq:runfreq}
\Prb(\mathrm{run})=
\begin{cases}
p,&0\le K<1/2,\\
p\rho,&1/2\le K<1,\\
0,&K=1.
\end{cases}
\end{equation}
A protected senior face \(1/2\) has price \(1/2\), and AT1 face \(1/4\)
has price \((1-p)/4\).
\end{proposition}
\begin{proof}
Below \(1/2\), both states run. Their \(A\) haircuts are respectively
\(K\) and zero, and each run returns half of \(B\)'s face to its
creditor. On \([1/2,1)\), the isolated state waits with
\((x,y)=(1,1-K)\), while the joint state runs with
\(x=K-1/2\). At \(K=1\) both states wait: \(A\)'s haircut is one
in each, while \(B\)'s are respectively zero and one half.
Weight these actual cash payments by \(p(1-\rho)\) and \(p\rho\)
and add full normal repayment with weight \(1-p\). This proves
\eqref{eq:QA}--\eqref{eq:QB}; it also shows why \(K=1\) must not
be filled in by a continuous interpolation of the \(A\) price.
The same run classification gives \eqref{eq:runfreq}. Senior debt
is always protected; AT1 is paid only outside stress.
\end{proof}

For comparison, the best conditional commitment at any fixed \(K\)
has no runs. Its designated prices are
\begin{equation}\label{eq:commitprices}
\begin{aligned}
Q_A^{\rm com}&=
\begin{cases}
1-p[K+(1-\rho)/2],&0\le K\le1/2,\\
1-p[(1-\rho)+\rho K],&1/2\le K\le1,
\end{cases}\\
Q_B^{\rm com}&=
\begin{cases}
\gamma(1-p/2),&0\le K\le1/2,\\
\gamma\{1-p[(1-\rho)(1-K)+\rho/2]\},&1/2\le K\le1 .
\end{cases}
\end{aligned}
\end{equation}
Indeed \eqref{eq:commitstates} gives \(A\) haircuts
\(K+1/2,K\) below \(1/2\), and \(1,K\) above it; \(B\)'s are
\(1/2,1/2\) below and \(1-K,1/2\) above. Taking the same state
expectations proves \eqref{eq:commitprices}; both pieces agree at
\(1/2\). Senior and AT1 prices are unchanged.

At \(p=1/2\), the discretionary values at the initially relevant
capital candidates are
\[
\begin{array}{c|cc}
K&Q_A&Q_B\\ \hline
0&1&3/8\\
1/2&(1+\rho)/2&3/8\\
1&1/2&1/2-\rho/8 .
\end{array}
\]
These are evaluations of the complete price functions, not a replacement
of the instrument interval by three admissible choices.

\subsection{Identification and its exact limitations}
If \(p\) is unknown, the maintained full AT1 stress loss and unit
risk-neutral discounting give \(p=1-4Q_{\rm AT1}\).
With an unknown common state-independent discount \(\delta\), the
protected senior benchmark gives \(\delta=2Q_{\rm senior}\);
divide the other prices by \(\delta\) before using their formulas.
State-dependent risk premia instead change physical probabilities
into state-price weights, as in Theorem~\ref{thm:pricepoly}. They
cannot be discarded when inferring physical stress frequencies or welfare.

\begin{proposition}[Sharp structural sets from the stated experiment]\label{prop:sharp}
Consider finitely many observed, known capital choices under common
\((p,\rho)\) and the declared pricing and equilibrium conventions.
Each discretionary observation is generated by the displayed discretionary
continuation. Each commitment observation is generated by the displayed
conditionally cost-minimizing schedules \eqref{eq:commitstates}; knowing
only that some irrevocable commitment exists would not suffice.
The sharp parameter set \(S\) is exactly the subset of \([0,1]^2\)
satisfying all applicable price, run and observed state-frequency equations
jointly. For that specified commitment experiment, the equations are
\eqref{eq:commitprices} and the zero-run restriction.

Fix the policy targets before using the observations: let
\(\kappa_{\rm D}(p,\rho)\) be the smallest cost-minimizing member of
\(\mathcal C(p,\rho)\) in Theorem~\ref{thm:capital}, and let
\(\kappa_{\rm C}(p,\rho)\) be the least member of the complete
commitment optimizer set in Theorem~\ref{thm:commitment}.
For either fixed target \(j\in\{\rm D,C\}\), its sharp identified
set is \(\{\kappa_j(p,\rho):(p,\rho)\in S\}\).
This targets a declared deterministic optimal capital choice, not every
optimal lottery or off-path commitment rule.
\end{proposition}
\begin{proof}
Necessity follows from the actual cash and run computations. Conversely,
every retained pair specifies the balance sheets and law above; the
proved authority and creditor rules generate every stipulated price
and behavior observation in that one model. Thus every retained
parameter and its policy target are attainable. The converse uses
common structural parameters across policies and claims, rather than
a separate recovery model for each observation.
\end{proof}

Under discretion, at \(K=0\), the two debt prices, the AT1 and senior prices, and
aggregate run frequency are all independent of \(\rho\). In particular,
run counts there do \emph{not} resolve the ambiguity left by prices.
At any \(K\in[1/2,1)\), known \(p>0\) and run frequency identify
\(\rho\); observed \(H\) labels among stress episodes also identify it.
The displayed price slopes supply the corresponding policy-variation
identification whenever their coefficient is nonzero. At \(p=0\),
these realized financial data do not identify the never-realized
conditional \(\rho\), although the hypothetical all-state credibility
tests remain specified.

Under discretion at \(p=1/2,K=0\), the same whole quoted claim-price vector is therefore
consistent with either optimal capital \(1/2\) or \(1\), depending on
which side of \(15/49\) the unidentified \(\rho\) lies. If the declared
target is worst-case incremental loss over the identified
\(\rho\in[0,1]\), Corollary~\ref{cor:robust} gives \(K=1\). This is
a specified robust counterfactual, not point identification of the
physical parameter.

\begin{proposition}[Identical prices and behavior, different credibility]\label{prop:crednonid}
If the authority's operating-cost coefficient is unknown, even the
complete prices, realized resolutions, run frequencies and observed
stress labels just described need not identify the all-history
credibility of a given announced rule.
\end{proposition}
\begin{proof}
For this comparison only, set \(K=0\) and replace \(x/4\) by \(cx\),
where \(c\in\{1/4,1/3\}\). Keep all other primitives and the common
pricing kernel fixed. Both economies are legal under the same cash
capacity: \(b_A+b_B\le5/4-x/2\), so even with the advance and
operating cost total disbursement is at most
\(7/4+(c-1/2)x\le7/4<2\).

After a run, for either \(H\ge0\), loss has \(x\)-slope \(c>0\),
so its unique optimum is zero. After waiting, first minimize \(x\)
at \([y-H]_+\). The reduced cost is
\[
 f_{H,c}(y)=1+H-y+c[y-H]_++y^2/2.
\]
Its derivative is negative through \(y=H\le1/2\), and thereafter
equals \(y+c-1\). The unique feasible minimum is
\[
 y=1-c>1/2,\qquad x=1-c-H .
\]
For both \(c\)'s and both states these actions satisfy all bounds.
The creditor therefore strictly runs in every stress state in both
economies. Actual resolutions have \(x=0\); all actual claim prices,
run frequencies and observed state labels coincide.

Announce the same complete rule in both economies: after running
choose \(x=0\); after waiting choose \(y=3/4\) and \(x=3/4-H\).
It is credible when \(c=1/4\). When \(c=1/3\) it is strictly
suboptimal at each unvisited waiting history, where the unique
alternative has \(y=2/3,x=2/3-H\). Investors anticipate that
deviation if waiting occurs; it still induces strict running, with
unchanged observed repayments. Thus the same observation law is
compatible with both credibility indicators.
\end{proof}

This comparison concerns unknown authority technology in an explicitly
specified two-model class. It does not apply the numerical
capital formula for \(c=1/4\) to a different \(c\). It shows directly
why expected loss allocation is not an observation of optimality at
unused legal histories. Known objective primitives or informative
counterfactual behavior restrictions are needed to make that inference.

\section{Full-solution scope, dependencies, and verification}\label{sec:scope}

\subsection{The scope declaration}
The exact target is \eqref{eq:original} for the fixed institution,
including its full legal rule set and both named institutional variants.
Theorem~\ref{thm:authority} proves every authority comparison and the
actual creditor response. Theorem~\ref{thm:capital} gives an exact optimum
for every \((p,\rho)\in[0,1]^2\); Theorem~\ref{thm:commitment} solves
commitment over the same entire capital interval and legal actions.
Theorem~\ref{thm:lottery} includes all the admitted lotteries rather
than using a deterministic restriction without proof. Both declared
failure states and all run/wait histories are checked, including those
assigned zero probability.

These results retain insured and senior claims, separate AT1 and
resolution debt, equity, actual cross-bank receipts, resource-constrained
recapitalization, statutory bridge operations, temporary funding and
endogenous runs simultaneously. The authority's conditional objective
and the ex ante welfare ledger are explicit; financial transfers and
real losses are counted separately. The nonlinear production loss is
derived from a stated constrained technology. Corollaries
\ref{cor:threshold} and \ref{cor:gap} exhibit an optimal instrument switch
and a strict commitment cost, rather than evaluating an exogenously
unique rule without an instrument.

Proposition~\ref{prop:prices} gives all-capital prices and actual run
frequencies; Proposition~\ref{prop:sharp} gives sharp structural sets
under the stated experiment. Proposition~\ref{prop:crednonid} explicitly
disproves identification of all-history credibility from the same
realized prices and behavior when authority costs are unknown. It
preserves the distinction between formal loss-imposition powers,
market expected repayments and sequential optimality.
Theorem~\ref{thm:pricepoly} states the general common-kernel restrictions
under which price bounds are sharp.

Theorems~\ref{thm:clearing}--\ref{thm:design} supply an additional exact
framework for the stated finite-history represented institutions.
It treats general legal continua, actual payment feedback, correct
belief consistency, pricing and budgets, and reports empty feasibility,
unbounded values and nonattainment separately. The concrete institution
proves attainment; the general theorem promises only the correct
endpoint or feasible \(\epsilon\)-optimal result. The framework is
not the whole claimed contribution by itself.

This answers the original fixed-system research task with its named
economic and identification requirements intact. It does not claim that
every banking system has this particular law or cost technology, that
general clearing is unique without leakage, that arbitrary infinite
strategic games admit the finite encoding, or that all creditor
indifference selections have the same optimal design. A different legal
action or information structure must be incorporated into its own full
domain. In the actual institution here none of its admitted continuous
authority actions or randomized rules remains unexamined.

The same scope test was applied to the prior attempts:
their quadratic exposure-cost and observable-risk theorems are preserved
and credited, while the clearing, run and nontrivial instrument problems
they expressly left unresolved are jointly solved here.
No broad worldwide research-priority claim follows from this comparison.
The original question requests explicit identification assumptions, not
execution of a particular empirical dataset. No replication of BIS data
or calibration to actual resolution agencies is claimed.

\subsection{External results and what was inspected}
The substantive external mathematical ingredients are effective real
quantifier elimination and algebraic sampling, as stated and used in
Theorem~\ref{thm:design}, and the finite-game Nash theorem used in
Lemma~\ref{lem:se}. The latter applies to a genuinely finite
perturbed agent game; the remainder of the continuation and consistency
argument is supplied here. The Kreps--Wilson citation credits the
sequential-equilibrium definition, whose exact consistency and
continuation tests were reproduced.

The primary Basu survey's quantifier-elimination statement, including
Theorem~2.27 in version~1, was inspected; the Basu--Pollack--Roy publisher
record was checked. Nash's primary article and Theorem~1 were inspected.
The Kreps--Wilson bibliographic record and abstract were checked; a
full-text mirror was unavailable during this authoring pass, and no
claim to inspecting all of that paper is made. The definitions,
encoding and fixed-prefix existence argument used here are explicit.
The Eisenberg--Noe primary article and publisher record supplied the
network context; our priority contraction and extreme-clearing proofs
do not import an unchecked uniqueness hypothesis from it.
Official BIS publication information and the bail-in abstract were
inspected. The catalogue's source audit identifies the economic
locators; neither empirical source is represented as proving this
strategic design theorem.

\subsection{Model, process, and actual checks}
The generating model is OpenAI GPT-6 based Codex. The system exposes that
generation, but not the exact backend variant or reasoning setting.
Existing same-session model-assisted reviewers inherited their model
settings; no different model was knowingly selected. Ji Ho Bae supplied
the task and constraints and is the submitter. No human mathematical
derivation, substantial mathematical correction, or independent human
proof verification was reported. Prompting, submission and compilation
are not treated as human mathematical authorship.

Research used the contribution policy's prompt~A and the exact canonical
statement and both manifest-listed priors pinned in the references.
Separate same-session reviewers examined the complete mathematical
modules and their original-scope combination. Revisions corrected
run-frequency identification, the full price curve, operating-cash
accounting, normal-state inclusion in pricing, and independence of a
public capital lottery from future shocks. The all-\(p\) boundary cases
and the lottery reduction were derived by hand. This is model-assisted
review, not independent human proof verification.

A separate remote exact-fraction corroboration record reports successful
finite checks: \(11{,}350\) ledger profiles, \(12{,}600\) authority
comparisons, \(4{,}900\) loss identities, \(4{,}900\) capital-candidate
comparisons, \(4{,}900\) price identities, \(194\) commitment flat-case
checks and the four counterexample states. These are finite arithmetic
checks only. The continuum, endpoint, equilibrium and global-optimality
claims rely on the proofs above, not on that enumeration. The script
identity and output are in the separate release record.

No general CAD computation, numerical global optimizer, empirical
estimation, Lean formalization or theorem-prover certificate was
performed for this proof. There is no independent human verification
and no exhaustive worldwide novelty certification. Separate reviewers
examined the whole manuscript and its specification corrections. Exact-version
review outcomes and executed typesetting/site checks are recorded separately
with actual source hashes. A successful build establishes neither
mathematical correctness nor novelty.

\section{Completion Estimate}
\noindent\textbf{COMPLETION: 100\%}
The full fixed-system task is claimed resolved: all legal continuation
comparisons, actual clearing and runs, the entire capital/probability
domain, both commitment variants, admitted lotteries, price-based sharp
sets and their nonidentification boundary are proved. The general
representation supplies the stated empty, unbounded and nonattained
cases and feasible epsilon-optimal rules. No mathematical gap within
this declared original-target scope is asserted; the percentage is a
scope/completion claim, not confidence or independent verification.

\begin{thebibliography}{99}
\raggedright
\bibitem{catalogue}
Economics problem catalogue, \emph{Credible bank resolution},
G28-66, JEL G28, source OP-0327. Canonical statement and common
conventions at repository commit
\texttt{b1f633121074e16579fa7fd05314e5f504d29bae}.
\href{https://github.com/mehmetmars7/Erdosproblems-llm-hunter/blob/b1f633121074e16579fa7fd05314e5f504d29bae/attacks/open_problems/economics/statements/G28-66.tex}{Pinned canonical statement}.
The welfare-optimal time-consistency formulation is identified there
as editorial and strategic.

\bibitem{bis}
A. Di Stefano, Y. Lengwiler and K. Rishabh (2026),
\emph{The credibility of bail-in}, BIS Working Paper 1356,
5 June 2026.
\href{https://www.bis.org/publications/working-paper-1356-credibility-bail-in}{Official publication and PDF}.
The catalogue locates the abstract and Section~7, printed p.~40
(PDF p.~42). Empirical source and motivation; not the source of our
optimal-policy theorem.

\bibitem{tucker}
P. Tucker (2019), \emph{Is the financial system sufficiently resilient:
a research programme and policy agenda}, BIS Working Paper 792.
Cover June 2019; release 4 July 2019.
\href{https://www.bis.org/publications/working-paper-792-financial-system-sufficiently-resilient-research-programme-and-policy-agenda}{Official publication and PDF}.
Catalogue locators: Section~5, printed pp.~12--16, and Section~7,
printed pp.~23--24. Policy-agenda context.

\bibitem{en}
L. Eisenberg and T. H. Noe (2001),
Systemic Risk in Financial Systems,
\emph{Management Science} 47(2), 236--249.
\href{https://doi.org/10.1287/mnsc.47.2.236.9835}{Publisher DOI};
\href{https://ms.mcmaster.ca/tom/Research\%20Papers/EiseNoe01.pdf}{Primary article copy}.

\bibitem{kw}
D. M. Kreps and R. Wilson (1982),
Sequential Equilibria,
\emph{Econometrica} 50(4), 863--894.
\href{https://doi.org/10.2307/1912767}{Publisher record};
\href{https://www.gsb.stanford.edu/faculty-research/publications/sequential-equilibrium}{Stanford author record}.

\bibitem{nash}
J. Nash (1951), Non-Cooperative Games,
\emph{Annals of Mathematics}, second series, 54(2), 286--295.
Theorem~1, p.~288.
\href{https://doi.org/10.2307/1969529}{Publisher DOI};
\href{https://irving.vassar.edu/faculty/gj/FTU/io/literature/Nash51.pdf}{Primary article copy}.

\bibitem{bpr}
S. Basu, R. Pollack and M.-F. Roy (2006),
\emph{Algorithms in Real Algebraic Geometry}, second edition,
Springer.
\href{https://doi.org/10.1007/3-540-33099-2}{Publisher record}.
Effective quantifier elimination and algebraic sampling.

\bibitem{basu}
S. Basu (2014), \emph{Algorithms in Real Algebraic Geometry: A Survey},
arXiv:1409.1534v1.
\href{https://arxiv.org/abs/1409.1534v1}{Versioned record};
\href{https://arxiv.org/pdf/1409.1534v1}{Versioned primary PDF}.
Theorem~2.27 is on printed p.~16; sampling is also derived above.

\bibitem{priorone}
GPT 6 Astra Ultra (2026),
\emph{Credible bank resolution: Systemic states, discretionary
write-downs, and risk selection}, first attempt, 9 October 2026.
\href{https://github.com/mehmetmars7/Erdosproblems-llm-hunter/blob/b1f633121074e16579fa7fd05314e5f504d29bae/attacks/open_problems/economics/gpt_6_astra_ultra/G28-66.tex}{Pinned original attempt}.
Its one-bank quadratic spillover and pricing results are credited.

\bibitem{priortwo}
GPT 6 Astra Ultra (2026),
\emph{Credible bank resolution: Systemic states, discretionary
write-downs, and risk selection}, version~2.
\href{https://github.com/mehmetmars7/Erdosproblems-llm-hunter/blob/b1f633121074e16579fa7fd05314e5f504d29bae/attacks/open_problems/economics/gpt_6_astra_ultra/G28-66_v2.tex}{Pinned second attempt}.
Its all-state quadratic exposure-cost and interacting-risk results
are credited; it expressly leaves clearing, runs and nontrivial
credible instrument design unresolved.
\end{thebibliography}
\end{document}
